\label{sec:background}

\subsection{Definition of Incumbent Pairing}

Let $\mathcal{I} = \{i_1, i_2, \ldots, i_n\}$ be the set of congressional incumbents in
a state who are seeking reelection following redistricting. Each incumbent $i_j$ has a
home census tract $t(i_j)$, determined from their FEC candidate filing address and
geocoded to the 2020 Census TIGER/Line tract layer. A redistricting plan $P$ assigns
each census tract to exactly one of $k$ congressional districts; let $d_P(t)$ denote
the district assignment of tract $t$ under plan $P$.

\begin{definition}[Incumbent Pairing]
Two incumbents $i_a$ and $i_b$ are \emph{paired} under plan $P$ if and only if their
home tracts are assigned to the same district:
\[
  d_P(t(i_a)) = d_P(t(i_b)).
\]
The \emph{pairing count} of plan $P$ is the number of unordered pairs of incumbents
that are paired:
\[
  C(P) = \left|\{(i_a, i_b) : i_a \neq i_b, d_P(t(i_a)) = d_P(t(i_b))\}\right|.
\]
The \emph{pairing rate} of plan $P$ is:
\[
  R(P) = \frac{C(P)}{\binom{n}{2}}.
\]
\end{definition}

Note that $R(P) = 0$ if and only if every district contains at most one incumbent's
home tract. $R(P) = 1$ if and only if all incumbents are assigned to the same district,
which is feasible only if $k=1$.

\subsection{Analytical Baseline Under Random Redistricting}

We derive the expected pairing count under a random redistricting process. The simplest
model assigns each incumbent independently and uniformly at random to one of $k$
districts. Under this model:

\begin{proposition}[Expected Pairing Count under Uniform Random Assignment]
If $n$ incumbents are assigned independently and uniformly at random to $k$ districts,
the expected number of paired pairs is:
\[
  E[C] = \binom{n}{2} \cdot \frac{1}{k}.
\]
\end{proposition}

\begin{proof}
For any pair $(i_a, i_b)$, $\Pr[d(i_a) = d(i_b)] = \sum_{j=1}^{k} \Pr[d(i_a)=j] \cdot
\Pr[d(i_b)=j] = \sum_{j=1}^k (1/k)^2 = 1/k$. Summing over all $\binom{n}{2}$ pairs
and applying linearity of expectation yields the result.
\end{proof}

This yields the expected pairing rates shown in Table~\ref{tab:baseline}. The variance
under this model is:
\[
  \text{Var}[C] = \binom{n}{2} \cdot \frac{1}{k} \cdot \left(1 - \frac{1}{k}\right)
  + 2\binom{n}{3} \cdot \frac{1}{k^2},
\]
where the second term arises from the covariance between pairs sharing a common incumbent.

\begin{table}[ht]
\centering
\caption{Analytical baseline for incumbent pairing under uniform random redistricting.}
\label{tab:baseline}
\begin{tabular}{l r r r r r}
\toprule
State & $k$ & $n$ & $E[C]$ & $\text{SD}[C]$ & $E[R]$ \\
\midrule
NC & 14 & 13 & 5.57 & 2.19 & 0.214 \\
WI & 8  & 8  & 3.50 & 1.62 & 0.250 \\
TX & 38 & 36 & 16.58 & 3.96 & 0.265 \\
\bottomrule
\end{tabular}
\end{table}

\subsection{Geographic Adjustment}

The uniform random model overstates the expected pairing rate because it ignores the
geographic adjacency structure of the redistricting problem. Incumbents whose home
tracts are geographically distant from each other are less likely to be paired by any
contiguous-district redistricting plan than the uniform model predicts. We compute a
geographic adjustment factor $\rho_s$ for each state $s$ by:

\begin{enumerate}
  \item Constructing the set of all contiguous-district plans consistent with the
    METIS adjacency graph structure.
  \item Computing, for each pair $(i_a, i_b)$, the fraction of randomly sampled plans
    that place both incumbents in the same district.
  \item Setting $\rho_s$ to the ratio of the geographically adjusted expected pairing
    count to the uniform-model expected pairing count.
\end{enumerate}

For NC, this adjustment yields $\rho_{\text{NC}} \approx 0.85$, reflecting the fact
that incumbent home tracts in NC are somewhat dispersed across the state's geography.
For WI, $\rho_{\text{WI}} \approx 0.70$, reflecting the relatively spread distribution
of WI incumbent residences. For TX, $\rho_{\text{TX}} \approx 0.78$.

Applying these adjustments, the geographically adjusted expected pairing rates are:
NC: 0.182; WI: 0.175; TX: 0.207. We use these adjusted expectations as the primary
baseline for testing.

\subsection{The Incumbency Protection Premium}

We define the \textit{incumbency protection premium} (IPP) of an enacted plan as the
difference between the geographically adjusted random baseline pairing rate and the
enacted pairing rate:
\[
  \text{IPP}(P) = E_{\text{geo}}[R] - R(P).
\]
A positive IPP indicates that the enacted plan provides more pairing protection than
random geography would imply; a zero or negative IPP indicates no deliberate protection.
Both enacted map analysis and algorithmic comparison can be expressed in terms of IPP.
